\chapter{Linear algebra: coordinates, maps, and modes}
\label{learn:linear_algebra}

\begin{learningcard}{Prerequisites and destination}
Bring precalculus vectors, dot products, matrix multiplication, and systems
of equations; use Chapter~\ref{learn:proof} for quantified statements.
We ask how a coupled problem can be expressed in coordinates that expose
independent directions. Eigenvectors will later become the modes of
Chapter~\ref{learn:dynamics}.
\end{learningcard}

\section{Vectors extend beyond arrows}
A real vector space is a set $V$ with addition and multiplication by real
scalars satisfying the following rules: addition is associative and
commutative, a zero vector and additive inverses exist, scalar multiplication
is associative with scalar products, one acts as the identity, and the two
distributive laws hold. The operations stay inside $V$.
The familiar space $\RR^n$ satisfies these rules component by component.
Polynomials of degree at most two form another example: adding two such
polynomials and scaling one preserves that degree bound.

A linear combination of vectors $v_1,\ldots,v_k$ is a sum
$\sum_{j=1}^k a_jv_j$ with real coefficients. Their span is the set of all
such combinations. They are linearly independent when
$\sum a_jv_j=0$ forces every $a_j=0$. A basis is an independent spanning
list. Every vector has unique coefficients in a basis: subtracting two
representations would give a vanishing linear combination, so independence
forces the coefficients to agree. The number of vectors in a basis of a
finite-dimensional space is its dimension.

For degree-at-most-two polynomials, $(1,x,x^2)$ is a basis. Coordinates
$(3,-2,5)$ represent $3-2x+5x^2$. Calling a polynomial a vector identifies
its algebraic operations; magnitude and spatial direction require additional
structure and interpretation.

\section{Worked example: a map becomes a matrix}
A map $T:V\to W$ is linear if
$T(au+bv)=aT(u)+bT(v)$ for all vectors and real scalars. Define
$T(a+bx+cx^2)=b+2cx$ on degree-at-most-two polynomials. This coefficient
rule defines a linear map using only algebra. Chapter~\ref{learn:calculus}
identifies the same rule as polynomial differentiation, $T(p)=p'$. Regard
the output as another polynomial of degree at most two. Since
$T(1)=0$, $T(x)=1$, and $T(x^2)=2x$, its matrix in the basis above is
\[
 A=\begin{pmatrix}0&1&0\\0&0&2\\0&0&0\end{pmatrix}.
\]
Each column is the coordinate vector of a basis image. Multiplication gives
$A(3,-2,5)^T=(-2,10,0)^T$, representing $-2+10x$ as required.
The kernel is the set of vectors mapped to zero, here the constant
polynomials. The image is the set of outputs, here polynomials of degree
at most one. Their dimensions are one and two.

The rank--nullity theorem states
$\dim\ker T+\dim\operatorname{im}T=\dim V$ for finite-dimensional $V$.
To see the mechanism, take a basis of the kernel and extend it to a basis
of $V$. The images of the added basis vectors span the image. Any linear
relation among those images would put a combination of the added vectors
inside the kernel, contradicting independence of the extended basis.
Thus the added images form a basis of the image.

\begin{learningcard}{Historical situation: an algebra of transformations}
Arthur Cayley's 1858 ``A Memoir on the Theory of Matrices'' treats matrices
as objects with operations and develops their algebra. The memoir is a
specific document in the nineteenth-century development of the subject;
methods for solving linear systems have much older histories. The Royal
Society publication context shows mathematical communication through a
learned institution, rather than a solitary discovery detached from readers
and prior techniques.\footnote{Arthur Cayley, ``A Memoir on the Theory of
Matrices,'' \emph{Philosophical Transactions of the Royal Society of London}
148 (1858), pp.~17--37, especially the opening definitions:
\url{https://doi.org/10.1098/rstl.1858.0002}.}
The polynomial example is a modern instructional construction.
\end{learningcard}

\section{Worked example: independent stretching directions}
An eigenvector of a square matrix $A$ is a nonzero vector $v$ satisfying
$Av=\lambda v$; its eigenvalue $\lambda$ is the stretching factor,
including a reversal when negative. For
\[
 A=\begin{pmatrix}2&1\\1&2\end{pmatrix},\qquad
 v_+=(1,1)^T,\qquad v_-=(1,-1)^T,
\]
direct multiplication gives $Av_+=3v_+$ and $Av_-=v_-$. These
independent vectors form a basis. Since
$(4,2)^T=3v_++v_-$, repeated application gives
$A^m(4,2)^T=3\cdot3^m v_++v_-$ for integers $m\geq0$.
A coupled recurrence becomes two scalar recurrences.

Normalize the vectors to obtain $q_+=v_+/\sqrt2$ and
$q_-=v_-/\sqrt2$. The matrix $Q$ with those columns obeys
$Q^TQ=I$, where $I$ is the identity matrix. Such a matrix is orthogonal;
its columns form an orthonormal basis, meaning unit lengths and pairwise
zero dot products. In these coordinates,
$Q^TAQ=\operatorname{diag}(3,1)$ and
$A=Q\operatorname{diag}(3,1)Q^T$.

\begin{figure}[htbp]
  \centering
  \begin{tikzpicture}[x=0.85cm,y=0.85cm,>=Stealth,font=\small,line width=0.8pt]
  \begin{scope}
    \draw[->] (-0.4,0)--(3.6,0) node[right] {$x$};
    \draw[->] (0,-1.5)--(0,3.6) node[above] {$y$};
    \node[above] at (1.7,4.1) {Input directions};
    \draw[->,line width=1.2pt] (0,0)--(1,1)
      node[above right] {$v_+=(1,1)^T$};
    \draw[->,dashed,line width=1.2pt] (0,0)--(1,-1)
      node[below right] {$v_-=(1,-1)^T$};
    \node[below left] at (0,0) {$0$};
    \foreach \value in {1,2,3} {
      \draw (\value,0.07)--(\value,-0.07);
      \draw (0.07,\value)--(-0.07,\value);
    }
  \end{scope}
  \begin{scope}[xshift=6.2cm]
    \draw[->] (-0.4,0)--(3.6,0) node[right] {$x$};
    \draw[->] (0,-1.5)--(0,3.6) node[above] {$y$};
    \node[above] at (1.7,4.1) {After applying $A$};
    \draw[->,line width=1.2pt] (0,0)--(3,3)
      node[above left] {$Av_+=3v_+$};
    \draw[->,dashed,line width=1.2pt] (0,0)--(1,-1)
      node[below right] {$Av_-=v_-$};
    \node[below left] at (0,0) {$0$};
    \foreach \value in {1,2,3} {
      \draw (\value,0.07)--(\value,-0.07) node[below] {$\value$};
      \draw (0.07,\value)--(-0.07,\value);
    }
  \end{scope}
  \node at (5.3,2) {$\displaystyle A=\begin{pmatrix}2&1\\1&2\end{pmatrix}$};
\end{tikzpicture}

  \caption{Read the input arrows at left, then the corresponding arrows at right on the same coordinate scale. The solid eigenvector grows by a factor of three; the dashed eigenvector keeps its length and direction. A general input splits into these two directions.}
  \label{fig:learn:linear_map}
\end{figure}

\section{Why symmetry guarantees enough eigenvectors}
The real symmetric spectral theorem states that every real matrix satisfying
$A^T=A$ has an orthonormal basis of real eigenvectors. Here is a proof
outline, with its analysis input made explicit. The continuous function
$q(v)=v^TAv$ attains a maximum on the unit sphere because that sphere is
compact in finite-dimensional Euclidean space. Let $u$ be a maximizer.
For any unit $w$ perpendicular to $u$, the curve
$u\cos t+w\sin t$ stays on the sphere. Differentiating $q$ along the
curve at $t=0$ gives $2w^TAu=0$. Consequently $Au$ has zero component
perpendicular to $u$, so $Au=\lambda u$.

For $w\perp u$, symmetry gives $u^TAw=(Au)^Tw=\lambda u^Tw=0$.
Thus the orthogonal complement of $u$ is preserved by $A$, and the
restriction there is again symmetric. Induction on dimension supplies an
orthonormal eigenbasis of that complement, which together with $u$ proves
the theorem. The compactness and derivative facts belong to real analysis;
the remainder is linear algebra.

A general real matrix need not have real eigenvectors: rotation by ninety
degrees sends every nonzero vector to a perpendicular vector. Even over
$\CC$, the matrix
\[J=\begin{pmatrix}1&1\\0&1\end{pmatrix}\]
has only a one-dimensional eigenspace and lacks an eigenbasis. These examples
explain why the theorem explicitly requires real symmetry.

\section{Exercises with full solutions}
\paragraph{1. Test a candidate vector space.}
Is $S=\{(x,y)\in\RR^2:x+y=1\}$ a subspace with ordinary operations?

\paragraph{Solution.}
A subspace must contain the zero vector, but $0+0\ne1$. Also $(1,0)$
belongs to $S$ while twice that vector does not. Therefore $S$ is an affine
line rather than a vector subspace. The parallel set $x+y=0$ is a subspace:
it equals the span of $(1,-1)$ and has dimension one.

\paragraph{2. Find a projection matrix and its modes.}
Find the orthogonal projection onto the line spanned by $(1,1)$.

\paragraph{Solution.}
A unit direction is $q=(1,1)^T/\sqrt2$. The projected vector is
$q(q^Tv)$, so
\[P=qq^T=\frac12\begin{pmatrix}1&1\\1&1\end{pmatrix}.\]
Multiplication gives $P(1,1)^T=(1,1)^T$ and $P(1,-1)^T=0$.
The eigenvalues are one along the line and zero perpendicular to it.
Applying the projection again preserves each component, so $P^2=P$.

\paragraph{3. Check an energy claim.}
For the symmetric example above, show $v^TAv>0$ for every nonzero real $v$.

\paragraph{Solution.}
Write $v=aq_++bq_-$. Orthonormality and the eigenvalue equations give
$v^TAv=3a^2+b^2$. A nonzero vector has at least one nonzero coefficient,
so the expression is positive. This property is called positive definiteness.
For a real symmetric matrix it holds exactly when every eigenvalue is
positive, since the same orthonormal expansion works in every dimension.
